\documentclass[10pt,a4paper]{amsart}
\usepackage[margin=19mm]{geometry}
\usepackage{amsmath,amssymb,mathtools}
\usepackage[scr=rsfs,cal=euler]{mathalfa}
\usepackage{xfrac}
\usepackage[unicode,hidelinks,bookmarks=false]{hyperref}
\newcommand{\ie}{i.e.\ }
\newcommand{\cf}{cf.\ }
\newcommand{\resp}{respectively\ }
\newcommand{\Subsection}{\subsection}
\providecommand{\nobreakdash}{\nobreak\hbox{-}}
\setlength{\emergencystretch}{2em}
\multlinegap0pt
\usepackage{bm}
\usepackage{bbm}
\usepackage{dsfont}
\usepackage{xspace}
\usepackage{cancel}
\usepackage{mathtools}
\usepackage{amsthm}
\makeatletter
\@ifundefined{theorem}{\newtheorem{theorem}{Theorem}}{}
\@ifundefined{claim}{\newtheorem{claim}[theorem]{Claim}}{}
\@ifundefined{proposition}{\newtheorem{proposition}[theorem]{Proposition}}{}
\makeatother
\newcommand{\R}{\mathbb{R}}
\newcommand{\Sphere}{\mathbb{S}^{n - 1}}
\newcommand{\Sset}{\mathscr S}
\newcommand{\esssup}{\operatorname*{ess\,sup}}
\title[Power-law asymptotics of fractional $L^p$ polar projection b]{Power-law asymptotics of fractional $L^p$ polar projection bodies}
\author{Tr\'i Minh L\^e}
\date{Source: arXiv:2601.05153v1 (CC BY 4.0)}
\begin{document}
\maketitle

\noindent\emph{Source and licence.} Power-law asymptotics of fractional $L^p$ polar projection bodies. Version arXiv:2601.05153v1 (CC BY 4.0), distributed under the Creative Commons Attribution 4.0 International licence, as recorded in the arXiv licence metadata for this exact version and shown on the abstract page https://arxiv.org/abs/2601.05153v1. Full rights notice: licenses/arxiv-2601.05153-CC-BY-4.0.txt.\par\smallskip

\noindent\textit{Editorial scope.} Counted unit: the Proposition whose source label is \texttt{prop.exist.infi.polar} in the article (source0.tex lines 510--525, source label \texttt{prop.exist.infi.polar}), delivered together with the article's own complete original proof of it (source0.tex lines 527--711, one \texttt{proof} environment of 185 lines). No mathematics is written, repaired or re-derived: statement and proof are the article's own text line for line. Required context retained uncounted: none. Every definition, display and label of those lines is retained unchanged. Omitted from this package and not redistributed: 1--509, 526--526, 712--2044. Nothing inside a retained excerpt is omitted or altered: every retained body is delivered exactly as the article's lines are written, so no omission receipt of any kind accompanies this package. Citations: the retained excerpts carry no citation command at all, so no bibliography entry of the article is transcribed and no reference is redistributed. Reference adaptation: none was needed and none was made; every reference inside the delivered unit resolves inside it, so no unresolved reference or citation is produced. Printed slips kept literal: no printed word, symbol, hypothesis or number of the counted statement or of its proof was altered. Layout and class: the article's own class is \texttt{article} with packages including amsmath, amsfonts, amssymb, mathtools, alphabeta, graphicx, amsthm, color, pgfplots, caption, subcaption, tikz, tikz-cd, float; the wrapper is the lane's pinned \texttt{amsart} editorial wrapper at 10pt on A4, which restates the theorem-like environments the retained text needs and adds the article's own macro definitions that the retained text needs (\texttt{\textbackslash R}, \texttt{\textbackslash Sphere}, \texttt{\textbackslash Sset}, \texttt{\textbackslash esssup}), with extra wrapper packages (bm, bbm, dsfont, xspace, cancel, mathtools). The delivered numbering is the unit's own. Assets: no retained span contains a figure environment, a graphics-inclusion command, a PDF-page inclusion, a table or a TikZ picture; no image file is copied into this package, the package declares no asset dependency and no image file is redistributed with it. Licence: version arXiv:2601.05153v1 of this article is distributed under the Creative Commons Attribution 4.0 International licence, as recorded in the arXiv licence metadata for this exact version and shown on the abstract page https://arxiv.org/abs/2601.05153v1. A full rights notice is delivered at licenses/arxiv-2601.05153-CC-BY-4.0.txt. Characters: every retained excerpt is pure ASCII, so the delivered engine, XeLaTeX with its default Latin Modern text font, reports no missing character.
\begin{proposition}\label{prop.exist.infi.polar}
    Let $f \in W^{1, 1}(\R^n) \cap W^{1, \infty}(\R^n)$ be nonzero.
    The following assertions hold true:
    \begin{itemize}
        \item[(i)] 
        For any fixed $s \in (0, 1)$, there exist $c = c(s, f) > 0$ and $\bar p = \bar p(s, f) > 1$ large enough such that $\Pi^{\ast, s}_p f \subset c B^n$ for every $p \geq \bar p$. 
        
        \item[(ii)] 
        For any fixed $s \in (0, 1)$, the set $\Pi^{\ast, s}_\infty \, f$ is origin-symmetric star body containing the origin in its interior.
        Furthermore, there exists $c = c(f) > 0$ and $\bar s = \bar s(f) \in (0, 1)$ such that $\Pi^{\ast, s}_\infty f \subset c B^n$ for every $s \in (\bar s, 1)$.

        \item[(iii)]
        The set $\Pi^{\ast}_\infty \, f$ is bounded origin-symmetric convex body with origin in its interior.
%        Furthermore, there exists $c = c(f, s) > 0$ such that $\Pi^\ast_\infty \, f \subset c B^n$.
    \end{itemize}
\end{proposition}

\begin{proof}
    \textit{(i)} 
    For any fixed $\alpha \in (0, 1)$, let us denote    
    \[
        \Sset_\alpha := \left\{ x \in \R^n : |f(x)| \geq \alpha \| f \|_{L^\infty(\R^n)} \right\}.
    \]
    
    \begin{claim}\label{claim.W_alpha_bounded}
    		For any fixed $\alpha \in (0, 1)$, the set $\Sset_\alpha$ is bounded.
    \end{claim}
	    
    \medskip
    
    \textit{Proof of Claim \ref{claim.W_alpha_bounded}}
    
	\medskip
	
    Set $L := \mathrm{Lip}(f)$ the Lipschitz constant of $f$ on $\R^n$, which is finite due to $f \in W^{1, \infty}(\R^n)$ and the Morrey embedding.
    Fix $\bar x \in \Sset_\alpha$ and $\bar r = \frac{\alpha \| f \|_{L^\infty(\R^n)}}{2L}$.
   The Lipschitz continuity of $f$ leads to
   	\begin{align*}
   		|f(x)| \geq |f(\bar x)| - L |x - \bar x| \geq \alpha \| f \|_{L^\infty (\R^n)} - L \bar r = \frac{\alpha}{2} \| f \|_{L^\infty(\R^n)}, \text{ for every } x \in B_{\bar r}(\bar x).
   	\end{align*}
	Thus,  $B_{\bar r}(\bar x) \subset \Sset_{\alpha/2}$.
    To prove the boundedness of $\Sset_\alpha$, arguing by contradiction, we assume that there exists a sequence $\{ x_k \} \subset \Sset_\alpha$ satisfying $|x_k| \to \infty$ as $k \to \infty$.
    We can further assume that $|x_k - x_{k'}| > 2\bar r$ and $B_{\bar r}(x_k) \cap B_{\bar r}(x_{k'}) = \varnothing$ for every $k, k' \in \mathbb{N}$.
	Due to the above observation, we note that $B_{\bar r}(x_k) \subset \Sset_{\alpha/2}$ for every $k$.
	It follows that
	\begin{align*}
		\int_{\R^n} |f(x)| dx \geq \sum_{k \in \mathbb N} \int_{B_{\bar r}(x_k)} |f(x)| dx \geq \frac{\alpha}{2} \| f \|_{L^\infty(\R^n)} \sum_{k \in \mathbb N} \underbrace{|B_{\bar r}(x_k)|}_{ \, = \, \bar r^n |B^n|} = \infty,
	\end{align*}
	which contradicts the fact that $f \in L^1(\R^n)$. \hfill $\lozenge$	
	
	\medskip
	
	We continue the proof of \textit{(i)}.
    Thanks to Claim \ref{claim.W_alpha_bounded}, the set $\Sset_{1/4}$ is bounded and so there exists $r_o > 0$ such that $\Sset_{3/4} \subset \Sset_{1/4} \subset r_o B^n$.
	Note that $0 < |\Sset_{3/4}| < + \infty$ and for any $p \in [1, + \infty)$, we have 
	\[
	\| f \|_{L^p(\Sset_{3/4})} \geq \frac{3}{4} \| f \|_{L^\infty(\R^n)} |\Sset_{3/4}|^{1/p}.
	\]
	Fix $\xi \in \Sphere$.
    For any $t > 2r_o$, the sets $r_oB^n$ and $\Sset_{3/4} - t\xi$ are disjoint and hence we get 
    \[
    	\| f \|_{L^p(\Sset_{3/4} - t \xi)} \leq \dfrac{1}{4} \| f \|_{L^\infty(\R^n)} |\Sset_{3/4} - t\xi|^{1/p} = \dfrac{1}{4} \| f \|_{L^\infty(\R^n)} |\Sset_{3/4}|^{1/p}.
    \]
    Using the triangle inequality, for any $t > 2r_o$, we have
    \begin{equation}\label{differ-Lp}
    	\begin{split}
    	\| f(\cdot + t \xi) - f \|_{L^p(\R^n)} 
    	\geq &  ~ 
    	~ \underbrace{\| f(\cdot + t \xi) - f \|_{L^p(\Sset_{3/4} - t\xi)}}_{= \, \| f   - f(~\cdot ~- t \xi) \|_{L^p(\Sset_{3/4})}} \\
    	\geq & ~ 
    	~ \| f \|_{L^p(\Sset_{3/4})} - \| f (\cdot - t \xi) \|_{L^p(\Sset_{3/4})} \\
    	\geq & ~
	\frac{3}{4} \| f \|_{L^\infty(\R^n)} |\Sset_{3/4}|^{1/p} - \dfrac{1}{4} \| f \|_{L^\infty(\R^n)} |\Sset_{3/4}|^{1/p} \\
    = & ~  \dfrac{1}{2}\| f \|_{L^\infty(\R^n)} |\Sset_{3/4}|^{1/p} > 0.
    	\end{split}
    \end{equation}
    Consequently, we obtain
    \begin{align*}
    	 \| \xi \|^{ps}_{\Pi^{\ast, s}_p \, f} 
         = & ~  p(1 - s) \int_0^\infty t^{-ps - 1} \| f(\cdot + t\xi) - f \|^p_{L^p(\R^n)} \, dt \\
    	 \geq & 
    	 ~ p(1 - s) \left( \dfrac{\| f \|_{L^\infty(\R^n)} |\Sset_{3/4}|^{1/p}}{2}  \right)^p \int_{2r_o}^\infty t^{-ps - 1}dt \\
    	 = & 
    	 ~ \dfrac{(1 - s)(2r_o)^{-ps}}{s}  \left( \dfrac{\| f \|_{L^\infty(\R^n)} |\Sset_{3/4}|^{1/p}}{2}  \right)^p,
    \end{align*}
   	and thus
   	\begin{align*}
   		\| \xi \|^{s}_{\Pi^{\ast, s}_p \, f} \geq \dfrac{1}{2} (2r_o)^{-s}\| f \|_{L^\infty(\R^n)} |\Sset_{3/4}|^{1/p} \left( \dfrac{1 - s}{s} \right)^{1/p}.
   	\end{align*}
 Since $\textstyle\lim_{p \to \infty} |\Sset_{3/4}|^{1/p}(s^{-1} - 1)^{1/p} = 1$, there exist $\bar p  = \bar p(s, f) > 1$ and $c = c(s, f) > 0$ such that $\| \xi \|_{\Pi^{\ast, s}_p \, f} \geq c$ for every $p \geq \bar p$. 
 Therefore, the conclusion follows.
   
    \medskip
    
    \textit{(ii)} 
%    We first observe that if $f \in W^{1, \infty}(\R^n) \cap W^{1, 1}(\R^n)$ is non-zero, then $\| \xi \|_{\Pi^{\ast, s}_\infty \, f} > 0$ for every $\xi \neq 0$ (see Appendix A). 
    Notice that 
    \[
        \| - \xi \|_{\Pi^{\ast, s}_\infty \, f} = \sup_{(x, t)} \dfrac{1}{t^s} |f(x - t\xi) - f(x)| = \sup_{(\tilde x, t)} \dfrac{1}{t^s} |f(\tilde x) - f(\tilde x + t\xi)| = \| \xi \|_{\Pi^{\ast, s}_\infty \, f},
    \]
    and hence $\Pi^{\ast, s}_\infty \, f$ is origin-symmetric.
    
    \medskip

    Let us check that the origin lies in the interior of $\Pi^{\ast, s}_\infty \, f$.
    One can directly see that $\xi \mapsto \| \xi \|_{\Pi^{\ast, s}_\infty \, f}^s$ satisfies the triangle inequality.
    Indeed, for each $\xi, \eta \in \R^n$, we have
    \begin{align*}
        \| \xi + \eta \|_{\Pi^{\ast, s}_\infty \, f}^s 
        = & ~ 
        \sup_{t > 0}
        \dfrac{1}{t^s} \| f(\cdot + t\xi + t\eta) - f \|_{L^\infty(\R^n)} \\
        \leq & ~
        \sup_{t > 0} 
        \dfrac{1}{t^s} \Big(
        \underbrace{\| f(\cdot + t\xi + t\eta) - f(\cdot + t \eta) \|_{L^\infty(\R^n)}}_{ = \, \| f(~\cdot~ + t\xi) - f \|_{L^\infty(\R^n)}} 
        + 
        \| f(\cdot + t\eta) - f \|_{L^\infty(\R^n)} 
        \Big) 
        \\
        \leq & ~ 
        \sup_{t > 0} \frac{1}{t^s} \| f(\cdot + t\xi) - f \|_{L^\infty(\R^n)} + \sup_{t > 0} \dfrac{1}{t^s}\| f(\cdot + t\eta) - f \|_{L^\infty(\R^n)} \\
        = & ~
        \| \xi \|^s_{\Pi^{\ast, s}_\infty \, f} 
        + \| \eta \|^s_{\Pi^{\ast, s}_\infty \, f} \,.
    \end{align*}
 	Notice that $\| t \xi \|_{\Pi^{\ast, s}_\infty \, f} = |t| \| \mathrm{sign}(t)\xi \|_{\Pi^{\ast, s}_\infty \, f}$ for every $t \in \R$ and $\xi \in \R^n$.
 	For any $x \in \R^n$, we write $x = \textstyle\sum_j x_j e_j$, where $\{ e_j \}_{j = 1}^d$ is the canonical basis of $\R^n$.
 	By the triangle inequality above, we get
 	\begin{equation}\label{origin-interior}
 		\| x \|_{\Pi^{\ast, s}_\infty f} \leq \left( \sum_j \| x_j e_j \|_{\Pi^{\ast, s}_\infty f}^s \right)^{1/s} = \left( \sum_j |x_j| \| \mathrm{sign}(x_j) e_j \|^s_{\Pi^{\ast, s}_\infty f}  \right)^{1/s} \leq a |x|,
 	\end{equation}
where $a > 0$ is independent of $x$. 
	This implies that $\Pi^{\ast, s}_\infty \, f$ contains the origin in its interior.

    \medskip
    
	To conclude that $\Pi^{\ast, s}_\infty \, f$ is a star body, it remains to prove that the gauge function $\| \cdot \|_{\Pi^{\ast, s}_\infty \, f}$ is continuous.
	Let $\{ \xi_k \}$ be a sequence converging to some $\xi \in \R^n$. 
 	Since the map $\xi \mapsto \| \xi \|_{\Pi^{\ast, s}_\infty \, f}^s$ satisfies the triangle inequality, we have
    \begin{align*}
        \Big\vert \| \xi_k \|_{\Pi^{\ast, s}_\infty \, f}^s - \| \xi \|_{\Pi^{\ast, s}_\infty \, f}^s \Big\vert
        \leq 
        \| \xi_k - \xi \|_{\Pi^{\ast, s}_\infty \, f}^s \underbrace{\, \leq \,}_{\eqref{origin-interior}} a^s | \xi_k - \xi|^s \to 0 \text{ as } k \to \infty.
    \end{align*}

    \medskip

    It remains to prove the boundedness of $\Pi^{\ast, s}_\infty \, f$ as $s$ is sufficiently close to $1$, that is, there exists $c = c(f) > 0$ such that $\| \xi \|_{\Pi^{\ast, s}_\infty \, f} \geq c$ for every $s$ closed to $1$.
    Let $r > 1$ be such that $\| f \|_{L^1(r B^n)} \geq \textstyle\frac{2}{3} \| f \|_{L^1(\R^n)}$.
    Fix $\xi \in \Sphere$.
    For any $t > 2r$, the two sets $rB^n$ and $rB^n - t\xi$ are disjoint and we also have $\| f \|_{L^1(rB^n - t\xi)} \leq \textstyle\frac{1}{3} \| f \|_{L^1(\R^n)}$.
    With these remarks, using the triangle inequality, we obtain the following estimate for any $t > 2r$,
    \begin{align*}
    	\| f(\cdot + t\xi) -  f\|_{L^\infty(\R^n)}
    	\geq & 
    	~ \| f(\cdot + t\xi) -  f\|_{L^\infty(rB^n - t\xi)} \\
    	\geq & 
    	~ \dfrac{1}{|rB^n - t\xi|} \underbrace{\| f(\cdot + t\xi)  - f \|_{L^1(rB^n - t\xi)} }_{\, = \, \| f - f(~\cdot~ - t \xi)\|_{L^1(rB^n)}} \\
    	\geq & 
    	~ \dfrac{1}{r^n |B^n|} \Big( \| f \|_{L^1(rB^n)} - \underbrace{\| f(\cdot - t\xi) \|_{L^1(rB^n)}}_{ \, = \, \| f \|_{L^1(rB^n - t\xi)} } \Big) 
     \\
    	\geq & 
    	~ \dfrac{1}{r^n |B^n|} \left( \dfrac{2}{3} \| f \|_{L^1(\R^n)} - \dfrac{1}{3} \| f \|_{L^1(\R^n)} \right) =  \dfrac{1}{3 r^n |B^n|} \| f \|_{L^1(\R^n)}.
    \end{align*}
    It follows from the above estimate and the definition of $\| \cdot \|_{\Pi^{\ast, s}_\infty \, f}$  that
	\begin{align*}
		\| \xi \|_{\Pi^{\ast, s}_\infty \, f} = \left( \sup_{t > 0} \dfrac{1}{t^s} \| f(\cdot + t\xi) - f \|_{L^\infty(\R^n)} \right)^{1/s} \geq \dfrac{\upsilon^{1/s}}{2r} 
		\quad \text{ with } \quad
		\upsilon = \dfrac{\| f \|_{L^1(\R^n)}}{3r^n |B^n|}.
	\end{align*}
	Note that $\textstyle\lim_{s \to 1^-} \upsilon^{1/s} = \upsilon$ and hence there exists $\bar s = \bar s(f) \in (0, 1)$ such that $\upsilon^{1/s} \geq \upsilon/2$ for all $s \in (\bar s, 1)$.
	This concludes the boundedness of $\Pi^{\ast, s}_\infty \, f$.
	
 	\medskip
 	
 	\textit{(iii)}
 	 The proof is direct.
 	 Indeed, we have, for every $\xi, \eta \in \R^n$, that
 	 \begin{align*}
 	 	\| \xi + \eta \|_{\Pi^\ast_\infty \, f} = & ~
        \esssup_{x \in \R^n} | \langle \nabla f(x), \xi + \eta \rangle| \\
        \leq & ~ \esssup_{x \in \R^n} | \langle \nabla f(x), \xi \rangle | + \esssup_{x \in \R^n} |\langle \nabla f(x), \eta \rangle| = \| \xi \|_{\Pi^\ast_\infty \, f} + \| \eta \|_{\Pi^\ast_\infty \, f}
 	 \end{align*}
 	 and
 	 \begin{align*}
 	 	\| \xi \|_{\Pi^\ast_\infty \, f } \leq \| \nabla f \|_{L^\infty(\R^n)} |\xi|.
 	 \end{align*}
 	 As a consequence of the above observations, we infer that $\xi \mapsto \| \xi \|_{\Pi^\ast_\infty \, f}$ is continuous.

        \medskip

        It remains to prove that $\Pi^\ast_\infty \, f$ is bounded, that is, $\textstyle\inf_{\xi \in \Sphere} \| \xi \|_{\Pi^\ast_\infty \, f} > 0$.
        Arguing by contradiction, we assume that $\textstyle\inf_{\xi \in \Sphere} \| \xi \|_{\Pi^\ast_\infty \, f} = 0$.
        Consequently, there exists $\xi_o \in \Sphere$ such that 
            \begin{align*}
                \esssup_{x \in \R^n} | \langle \nabla f(x), \xi_o \rangle | = 0.
            \end{align*}
Consequently, for a.e $x \in \R^n$, $f$ is constant along the line $x + t\xi_o$.
        Since $f \in L^1(\R^n)$, we infer that $f \equiv 0$, which makes a contradiction.
 	The proof of Proposition \ref{prop.exist.infi.polar} is complete.
\end{proof}

\end{document}
